\section{Latent diffusion：把贵的部分压缩掉}

\subsection{为什么要进 latent space}

把 diffusion 直接跑在高分辨率像素上太贵了。空间分辨率一高，token 数量和计算成本就会迅速膨胀。于是一个很实用的策略出现了：先用 encoder 把图像压缩到更小的 latent，再在 latent 上做 diffusion，最后再用 decoder 还原。

\begin{figure}[H]
\centering
\includegraphics[width=0.95\textwidth]{slides-images/slide-084.jpg}
\caption{直接在高分辨率图像上做 diffusion 不够经济，latent diffusion 的动机正是压缩空间分辨率\protect\footnotemark}
\end{figure}
\footnotetext{Slides 第85页。}

\begin{figure}[H]
\centering
\includegraphics[width=0.95\textwidth]{slides-images/slide-085.jpg}
\caption{Latent diffusion 的结构：图像先被编码成 latent，再在 latent 上做 diffusion，再解码回图像\protect\footnotemark}
\end{figure}
\footnotetext{Slides 第86页。}

\begin{importantbox}{latent diffusion 为什么成为工业主流}
这一步不是“稍微省点显存”这么简单，而是把生成难题转移到了更适合建模的空间里。latent 维度更低，模型更容易学到语义结构，采样成本也明显下降。
\end{importantbox}

\subsection{VAE 是 latent diffusion 的前半段}

latent diffusion 里 encoder / decoder 通常来自 VAE。训练时需要同时优化重建项和 KL prior，但 KL 往往会被压得很小，因为这里的目标是压缩和保真，而不是强约束 latent 服从某种过强的先验。

\begin{figure}[H]
\centering
\includegraphics[width=0.95\textwidth]{slides-images/slide-093.jpg}
\caption{如何训练 encoder 和 decoder：本质上就是一个 VAE\protect\footnotemark}
\end{figure}
\footnotetext{Slides 第94页。}

\begin{figure}[H]
\centering
\includegraphics[width=0.95\textwidth]{slides-images/slide-094.jpg}
\caption{VAE 的现实问题：decoder 往往会发糊，单靠重建项很难拿到很锐利的细节\protect\footnotemark}
\end{figure}
\footnotetext{Slides 第95页。}

\begin{warningbox}{为什么纯 VAE 容易模糊}
VAE 的 decoder 追求的是条件平均意义上的重建最优，而不是感知质量最优。对于多解的高频细节，平均往往意味着“糊”。这就是为什么后续工程里经常要引入 GAN 式的对抗约束来补锐度。
\end{warningbox}

\subsection{为什么现代系统是 VAE + GAN + diffusion}

当 decoder 细节不够锐利时，就可以加一个 discriminator 帮它把输出推向更自然的图像外观。于是 latent diffusion 的前半段是 VAE，后半段是 diffusion，局部细节修正则借用了 GAN 的思路。

\begin{figure}[H]
\centering
\includegraphics[width=0.95\textwidth]{slides-images/slide-096.jpg}
\caption{一个真实的现代管线不是“单模型”，而是 VAE + GAN + diffusion 的组合\protect\footnotemark}
\end{figure}
\footnotetext{Slides 第97页。}

\begin{knowledgebox}{现代 LDM 管线的分工}
\begin{itemize}
    \item VAE 负责压缩和还原
    \item GAN 负责局部质感和清晰度
    \item diffusion 负责全局生成和条件控制
\end{itemize}
\end{knowledgebox}

\begin{figure}[H]
\centering
\includegraphics[width=0.95\textwidth]{slides-images/slide-091.jpg}
\caption{Latent diffusion 的基本工作流：先采样随机 latent，再逐步去噪，最后经 decoder 输出图像\protect\footnotemark}
\end{figure}
\footnotetext{Slides 第92页。}

